% Generated draft: requires mathematical and source review.
\documentclass[11pt]{article}
\usepackage[margin=25mm]{geometry}
\usepackage{fontspec}
\setmainfont{Latin Modern Roman}
\usepackage{amsmath,amssymb,mathtools,unicode-math}
\setmathfont{Latin Modern Math}
\usepackage{enumitem,needspace,xurl,hyperref}
\setlength{\parindent}{0pt}
\setlength{\parskip}{5pt}
\setlength{\emergencystretch}{3em}
\newcommand{\N}{\mathbb{N}}
\newcommand{\Z}{\mathbb{Z}}
\newcommand{\Q}{\mathbb{Q}}
\newcommand{\R}{\mathbb{R}}
\newcommand{\C}{\mathbb{C}}
\newcommand{\F}{\mathbb{F}}
\newcommand{\A}{\mathbb{A}}
\newcommand{\PP}{\mathbb{P}}
\newcommand{\E}{\mathbb{E}}
\newcommand{\eps}{\varepsilon}
\newcommand{\dd}{\,\mathrm{d}}
\begin{document}
\section*{UnsolvedMath ID: 30004367}
\textbf{Source problem code:} catalogue entry
\subsection{Definitions and mathematical statement}
Let \(\mathcal T\) denote the ambient triangulated category mentioned in the supplied record, and let \(\mathcal H\) denote the category called its heart. A heart normally means the abelian category \(\mathcal T^{\leq 0}\cap\mathcal T^{\geq 0}\) of a specified \(t\)-structure on a triangulated category. The source does not specify that \(t\)-structure or unambiguously identify the category to which the phrase ``its heart'' refers. Write \(D^b(\mathcal H)\) for the bounded derived category of \(\mathcal H\), obtained from bounded complexes by inverting quasi-isomorphisms, and write the supplied realization functor as
\[
G\colon D^b(\mathcal H)\longrightarrow\mathcal T.
\]
For a functor \(G\), its essential image consists of the objects isomorphic to some \(G(X)\). Thus the source's assumption that the abelian category \(\mathcal A\) is contained in the essential image means, at the level of objects,
\[
\forall A\in\operatorname{Ob}(\mathcal A)\quad
\exists X\in D^b(\mathcal H)\quad A\cong G(X)\text{ in }\mathcal T.
\]
This presupposes a way to regard objects of \(\mathcal A\) as objects of \(\mathcal T\); the embedding and its compatibility with the abelian and triangulated structures are not specified.

\textbf{Question as supplied.} Under the intended realization-functor setup, does
\[
\mathcal A\subseteq\operatorname{EssIm}(G)
\quad\Longrightarrow\quad
G\colon D^b(\mathcal H)\xrightarrow{\sim}\mathcal T
\]
hold? Here the conclusion is interpreted as saying that \(G\) is an equivalence of triangulated categories: it is a triangle functor, every object of \(\mathcal T\) is isomorphic to some \(G(X)\), and for every \(X,Y\in D^b(\mathcal H)\) the induced map
\[
\operatorname{Hom}_{D^b(\mathcal H)}(X,Y)
\longrightarrow
\operatorname{Hom}_{\mathcal T}(G(X),G(Y))
\]
is bijective.

\textbf{Unresolved scope.} The record does not identify \(\mathcal T\) with \(D^b(\mathcal A)\), state how \(\mathcal A\) relates to \(\mathcal H\), impose a generation condition on \(\mathcal A\), or give the hypotheses and construction ensuring that the realization functor exists. These omissions prevent an unambiguous, fully quantified formulation of the intended general problem; the displayed implication records the supplied target with those qualifications unresolved.

The supplied dated assessment reports an affirmative hereditary case, where an abelian category is called hereditary when its higher extension groups \(\operatorname{Ext}^n(A,B)\) vanish for all objects \(A,B\) and all \(n\geq 2\). The record does not specify the full realization setup for this reported case. Hereditariness is not imposed in the general question, and this reported conditional result is not independently verified here.
\subsection{Short English statement}
The question asks whether a realization functor from the bounded derived category of a heart must be a triangulated equivalence when its essential image contains the given abelian category \(\mathcal A\). The ambient category, heart, and their relationship to \(\mathcal A\) need clarification before the intended statement is fully precise. The supplied status is \emph{partially solved}, with a reported hereditary case; this is dated curation rather than an independently verified conclusion.
\subsection{Sources}
\begin{itemize}
\item Supplied local UnsolvedMath registry record, permanent ID \(30004367\), title \emph{Realization Functors as Derived Equivalences}, exact source problem code \texttt{catalogue entry}, in the Oberwolfach Reports Open Problems collection. The record gives proposed year 2020 and names no individual proposer.
\item Supplied source citation: \emph{Representation Theory of Quivers and Finite Dimensional Algebras} (2021), Oberwolfach Reports, pp.~143--230, DOI \texttt{10.4171/owr/2020/3}. \url{https://doi.org/10.4171/owr/2020/3}. The supplied background describes Oberwolfach Reports as documenting mathematical workshops and their open-problem discussions.
\item Supplied literature reference: Chen--Han--Zhou, \emph{Derived equivalences via HRS tilting}. \url{https://doi.org/10.1016/j.jpaa.2019.106749}.
\item Supplied literature reference: \emph{Oberwolfach report, realization functors}. \url{https://ems.press/content/serial-article-files/46842}.
\item Literature assessment dated exactly \texttt{2026-08-22}: supplied status \emph{partially solved}. The assessment states that the realization functor is an equivalence when the abelian category is hereditary, that the question extends beyond that conditional case, and that it ``has no general affirmative solution.'' This is supplied dated curation derived from the OWR record, not independently verified literature assessment. The supplied background requests independent verification of time-sensitive status and cited claims.
\item The registry labels the statement \texttt{corrected\_verified} and reports manual semantic, typographic, and LaTeX review asserting self-containment and preservation of the mathematical target. Those are supplied editorial claims; they do not resolve the missing setup identified in this draft.
\end{itemize}
\end{document}
